\spread{Learning from every way to write the same ink}
\noindent A name may be written one glyph at a time, copied as a phrase, or copied from more than one position. We trained a writer to credit each way of producing the human answer and tested which facts survived in its own reply.
\platehere{Figure_166_learning_complete_answers}{Two answers with the evidence region held fixed across writers}{Each panel shows the original pixels for a pictured question, the human reference and two freely generated replies. The panels are the first gains in either direction on the second English confirmation set. Panel a recovers the object “public support”. Panel b retains the date but loses the restriction “online”, changing the claim about when the publication began.}
\begin{multicols}{2}
\section*{A probability for the complete image}
Let $w_{t,k}$ sum the probabilities of the actions that write the next $k$ reference glyphs at position $t$. Subsequent state depends on the emitted ink. Starting with $A_0=1$, compute $A_j=\sum_{t+k=j}A_t w_{t,k}$ in reading order and multiply the final mass by the stop probability. Training maximizes this sum over segmentations. Writing uses the same greedy action decoder as the longest-run control. The reverse recurrence gives the posterior mass of each compatible edge. Lean verifies a finite edge-weight perturbation identity. On three actual English and Chinese training examples the full parameter gradient agreed with an independent differentiable reference to within $4.6\times10^{-8}$ in relative Euclidean norm.
\section*{Which facts survive?}
We reserved another 96 English and 96 Chinese passages before training and drew two questions per passage. In English the single-glyph, longest-run and all-path writers gave 11, seven and ten exact images out of 192. Among the 188 answers absent from the input their counts were 11, five and ten. The requested fact appeared in 43, 48 and 33 replies. Exact answers to both questions occurred on zero, one and one passage. The three writers all gave zero exact images on the Chinese reading questions. Fact presence fell from 25 for longest-run training to five for all paths.
\section*{Probability and meaning}
On the 192 English references the all-path writer averaged 0.339 nats per glyph, compared with 0.341 for single-glyph training. Among 176 different stopped replies it gave the human reference greater probability than its own output only 19 times. A blinded judge recognized all 16 original English human references, but only nine after a fact substitution. Paired semantic generalization remains unresolved.
\end{multicols}
\newpage
\spread{When does the writer learn to state the fact?}
\noindent We followed the all-path model during fitting, then held the pictured question and fact fixed and changed which writer joined them into a sentence.
\platehere{Figure_167_learning_from_pictured_answers}{Training trajectories on previously held passage families}{Each marker is measured from a saved model on the same 192 questions per language. Fact presence requires the requested short answer in a stopped output. Exactness requires the entire ordered bitmap sequence and stop. The dotted line marks the missing 2,400-update checkpoint. The independent confirmation set on the preceding page was reserved for the final weights.}
\begin{multicols}{2}
\section*{What more updates bought}
English exact images increased from zero at initialization to eight at the final checkpoint. Fact presence reached 41 at 4,800 updates and finished at 36. On Chinese reading questions the final writer stopped on 191 of 192 and matched zero complete references. On 35 Chinese catalogue questions it progressed from zero to eight exact images. The two run writers each saw the same 89,263 full human responses four times. The longest-run fit took 180 seconds and 22.2 kJ on one RTX 4090 D; all-path fitting took 190 seconds and 22.9 kJ. Both reuse the same selector (70.5 seconds, 8.1 kJ). Preparation, evaluation and host energy are excluded.
\section*{Supplying the missing information}
On the independent 192 English questions, giving the all-path writer the human-marked region raised complete exact replies from ten to 11. Replacing the entire region with the human short answer yielded five. With the same pictured question and short answer, an earlier writer explicitly trained to formulate a sentence from a fact yielded 34. It matched both the original and edited sentence on ten of 43 substitutions, against three for the all-path writer. Swapping sibling questions while keeping the short fact fixed eliminated all 34 exact matches to the original answers. A supplied human fact changes the input distribution; the older specialist also has a different training history.
\section*{An independent reading benchmark}
On 16 matched image prompts the local quantized Qwen3-VL-8B included the requested fact on six of eight English and six of eight Chinese questions, versus zero and one for the all-path writer. Inference used 4.86 kJ and 0.107 kJ of measured GPU energy respectively. Their quality differed on these inputs. The all-path writer failed all 152 eight-glyph book continuations and 260 historical specimen questions. An independently predicted fact image and a writer trained on such predictions remain to be tested together against an equally trained direct writer.
\end{multicols}
\newpage
